\chapter{Gestion de version}\label{chap:git}

\section{Principes}
Le dépôt \code{jee-ecommerce-jsf-jpa} est public sur GitHub. Il suit les règles de l'atelier~2 : commits atomiques qui
compilent, messages au format \textit{Conventional Commits} \cite{conventional}, et aucun secret versionné
(\code{db.properties} est ignoré ; seul \code{db.properties.example} est suivi).

\section{Historique}
Le tableau~\ref{tab:commits} donne la séquence. Chaque commit correspond à une couche ; le rapport est ajouté en dernier.

\begin{table}[htbp]
\centering
\caption{Séquence de commits}\label{tab:commits}
\small
\begin{tabularx}{\linewidth}{@{}lX@{}}
\toprule
\textbf{Étape} & \textbf{Message}\\
\midrule
Projet & \code{chore: initialise Maven project, database script and configuration template}\\
Persistance & \code{feat(persistence): add JPA entities, utilities and schema tool}\\
Services & \code{feat(service): add business services and JPA implementations}\\
Web & \code{feat(web): add JSF managed beans, filters and configuration}\\
Vues & \code{feat(views): add Facelets layout, composite components and pages}\\
Tests & \code{test: add business, security and Facelets static checks}\\
Rapport & \code{docs(report): add LaTeX report, figures and PDF}\\
\bottomrule
\end{tabularx}
\end{table}
